\documentclass[12pt]{article}
\usepackage[utf8]{inputenc}
\usepackage[T1]{fontenc}
\usepackage{hyperref}
\title{Tartessian Language Resources for K-12 Curriculum Development}
\author{Compiled by Claude Sonnet 5 (claude-sonnet-5) for Andr\'as Kornai}
\date{}
\begin{document}
\maketitle
\section{Introduction}
Tartessian (ISO 639-3 txr, Glottocode tart1237) is an extinct language of the south-west Iberian Peninsula, known only from roughly 75 to 100 inscriptions (counts vary by source) in the Southwestern script, usually dated from the 7th century BC. Its classification is disputed; Koch argues it is Celtic, but this is not settled. There are no speakers, so a school curriculum in the language is not possible; resources are scholarly only.
\section{Primary Educational Resources}
None exist in search results.
\section{Linguistic and Reference Materials}
\begin{itemize}
\item Wikipedia, Tartessian language: \url{https://en.wikipedia.org/wiki/Tartessian_language}
\item Koch, on the classification of the SW inscriptions: \url{https://aemap.ac.uk/media/uploads/PDFs/sw_classification_koch.pdf}
\item Koch, Palaeohispanica 10: \url{https://aemap.ac.uk/media/uploads/PDFs/koch_palaeohispanica_10.pdf}
\item University of Barcelona repository document: \url{https://diposit.ub.edu/dspace/bitstream/2445/217425/1/188238.pdf}
\end{itemize}
\section{Religious and Community Resources}
None found.
\section{Digital and Audio Resources}
\begin{itemize}
\item Unicode proposal for Palaeohispanic scripts: \url{https://www.unicode.org/L2/L2015/15119-palaeohispanic.pdf}
\item Omniglot, Iberian scripts: \url{https://omniglot.com/writing/iberian.htm}
\end{itemize}
\section{Curriculum Development Context}
At most the inscriptions could support enrichment units on ancient Iberian history and writing systems.
\end{document}